% standard beamer lecture template for slides
% by Derek Huang
\documentclass{beamer}
\usepackage[utf8]{inputenc}
\usepackage{algorithm2e, amsmath, amssymb, amsfonts, graphicx}
% allow section.equation numbering
\numberwithin{equation}{section}
% use boadilla theme
\usetheme{Boadilla}
% remove navigation symbols
\usenavigationsymbolstemplate{}
% get numbered figure captions
\setbeamertemplate{caption}[numbered]
% changes itemize to circle + other things
\useoutertheme{split}
\useinnertheme{circles}

% command for the title string. change for each lecture
\newcommand{\lecturetitle}{Linear Regression}
% allow automatic alert-highlighted references, href
\newcommand{\aref}[1]{\alert{\ref{#1}}}
\newcommand{\ahref}[2]{\href{#1}{\alert{#2}}}
% title page stuff. brackets content displayed in footer bar
\title[\lecturetitle]{\lecturetitle}
% metadata. content in brackets is displayed in footer bar
\author[Derek Huang (BAC Advanced Team)]{Derek Huang}
\institute{BAC Advanced Team}
\date{March 4, 2021}

% change "ball" bullet to numbered bullet and section title for section
\setbeamertemplate{section in toc}{\inserttocsectionnumber.~\inserttocsection}
% change ball to gray square (copied from stackoverflow; \par needed for break)
\setbeamertemplate{subsection in toc}{        
    \hspace{1.2em}{\color{gray}\rule[0.3ex]{3pt}{3pt}}~\inserttocsubsection\par}
% use default enumeration scheme
\setbeamertemplate{enumerate items}[default]
% required line that fixes the problem of \mathbf, \bf not working in beamer
% for later (post-2019) TeX Live installations. see the issue on GitHub:
% https://github.com/josephwright/beamer/issues/630
\DeclareFontShape{OT1}{cmss}{b}{n}{<->ssub * cmss/bx/n}{}

\begin{document}

% title slide
\begin{frame}
    \titlepage
    \centering
    % relative path may need to be updated depending on .tex file location
    \includegraphics[scale = 0.1]{../bac_logo1.png}
\end{frame}

% table of contents slide
\begin{frame}{Overview}
    \tableofcontents
\end{frame}

\section{Linear regression}

\begin{frame}{Motivation}
    \begin{itemize}
        \item
        The simplest\footnote{
            It may seem simple, but it contains more subtleties than one
            might expect.
        } regression method that one should understand.

        \item
        \textit{Remark.} Simple $ \ne $ not effective.

        \item
        The fundamental basis for understanding more ``interesting'' models.

        \item
        Methods used to fit the model in practice are less expensive in both
        CPU time and memory usage compared to other models.

        \item
        How much about the model do you \textit{really} understand?
    \end{itemize}
\end{frame}

\subsection{The linear model}

\begin{frame}{The linear model}
    \begin{itemize}
        \item
        Suppose on the real-world $ (\Omega, \mathcal{F}, \mathbb{P}) $ we
        have random variables $ X : \Omega \rightarrow \mathbb{R}^d $,
        $ Y : \Omega \rightarrow \mathbb{R} $. The \textit{linear regression
        model}\footnote{
            Technically, it should be called \textit{affine} regression since
            $ b $ provides a shift.
        } posits that
        \begin{equation} \label{lr_model}
            Y = \mathbf{w}^\top X + b + \varepsilon
        \end{equation}
        Here $ \mathbf{w} \in \mathbb{R}^d $, $ b \in \mathbb{R} $, and
        independent $ \varepsilon \sim \mathcal{N}(0, \sigma^2) $\footnote{
            The assumption on homoscedastic Gaussian errors will be discussed
            next lecture.
        }, $ \sigma \in (0, \infty) $.

        \item
        Interpretation is that $ Y $ can be written as a translated linear
        combination of the $ X $ components plus an independent noise term.

        \item
        The perennial question: Assuming true $ \mathbf{w}, b $ exists, how do
        we estimate $ \mathbf{w}, b $ sufficiently well given training data
        $ \mathcal{D} $?
    \end{itemize}
\end{frame}

\subsection{Solving the system}

\begin{frame}{Solving the system}
    \begin{itemize}
        \item
        Suppose $ \mathcal{D} \triangleq (\mathbf{X}, \mathbf{y}) $, where
        $ \mathbf{X} \in \mathbb{R}^{N \times d} $ is the input matrix, each
        row a transposed data example, $ \mathbf{y} \in \mathbb{R}^N $ the
        response vector. Clearly the problem we have is that of finding
        $ \hat{\mathbf{w}} \in \mathbb{R}^d , \hat{b} \in \mathbb{R} $
        such that
        \begin{equation*}
            \mathbf{X}_\mathbf{1}\hat{\mathbf{w}}_{\hat{b}} \triangleq
            \begin{bmatrix}
                \ \mathbf{X} & \mathbf{1} \
            \end{bmatrix}
            \begin{bmatrix}
                \ \hat{\mathbf{w}} \ \\ \ \hat{b} \
            \end{bmatrix} \approx \mathbf{y}
        \end{equation*}
        Here $ \mathbf{X}_\mathbf{1} \in \mathbb{R}^{N \times (d + 1)} $,
        $ \hat{\mathbf{w}}_{\hat{b}} \in \mathbb{R}^{d + 1} $. Assume
        $ \operatorname{rk}(\mathbf{X}_\mathbf{1}) = \min\{N, d + 1\} $.

        \item
        Problem: often $ N \ne d + 1 $, so $ \mathbf{X}_\mathbf{1} $ is
        \textbf{not} invertible, and therefore $ \nexists
        \hat{\mathbf{w}}_{\hat{b}} $ such that $ \hat{\mathbf{w}}_{\hat{b}} =
        \mathbf{X}_\mathbf{1}^{-1}\mathbf{y} $, i.e. $ \mathbf{X}_1
        \hat{\mathbf{w}}_{\hat{b}} = \mathbf{y} $. But consider
        \begin{equation} \label{lr_obj_aug}
            \begin{array}{ll}
                \displaystyle\min_{\mathbf{w}_b} &
                \Vert\mathbf{y} - \mathbf{X}_\mathbf{1}\mathbf{w}_b\Vert_2^2
            \end{array}
        \end{equation}
    \end{itemize}
\end{frame}

\begin{frame}{Solving the system}
    \begin{itemize}
        \item
        Cast solving of an over/under-determined linear system as minimization
        of the $ \ell^2 $-norm of the residual vector $ \mathbf{r} \triangleq
        \mathbf{y} - \mathbf{X}_\mathbf{1}\hat{\mathbf{w}}_{\hat{b}} $.

        \item
        \textit{Remark.} (\aref{lr_obj_aug}) minimizes the \textbf{squared}
        $ \ell^2 $-norm of $ \mathbf{r} $, but note that
        \begin{equation*}
            \arg\min_{\mathbf{w}_b}
            \Vert\mathbf{y} - \mathbf{X}_\mathbf{1}\mathbf{w}_b\Vert_2 =
            \arg\min_{\mathbf{w}_b}
            \Vert\mathbf{y} - \mathbf{X}_\mathbf{1}\mathbf{w}_b\Vert_2^2
        \end{equation*}
        For function $ f $, $ \arg\min_\mathbf{x}f(\mathbf{x}) $ retrieves
        the point that minimizes $ f $.

        \item
        \textit{Remark.} (\aref{lr_obj_aug}) minimizes the
        \textbf{sum of squared residuals}. Also, note that
        $ \Vert\mathbf{y} - \mathbf{X}_\mathbf{1}\mathbf{w}_b\Vert_2^2 =
        (\mathbf{y} - \mathbf{X}_\mathbf{1}\mathbf{w}_b)^\top(\mathbf{y} -
        \mathbf{X}_\mathbf{1}\mathbf{w}_b) $ by definition.

        \item
        \textit{Properties of }(\aref{lr_obj_aug}). Let $ f(\mathbf{w}_b) =
        \Vert\mathbf{y} - \mathbf{X}_\mathbf{1}\mathbf{w}_b\Vert_2^2 $.
        \begin{enumerate}
            \item
            $ \forall \mathbf{w}_b \in \mathbb{R}^{d + 1} $,
            $ \exists \nabla f(\mathbf{w}_b) $ (differentiable everywhere).

            \item
            $ f $ convex $ \Rightarrow \hat{\mathbf{w}}_{\hat{b}} $ optimal
            $ \Leftrightarrow \nabla f(\hat{\mathbf{w}}_{\hat{b}}) =
            \mathbf{0} $ ($ \exists $ global minimum) \cite{bv_convex_opt}.
        \end{enumerate}
    \end{itemize}
\end{frame}

\begin{frame}{Solving the system}
    \begin{itemize}
        \item
        At optimal $ \hat{\mathbf{w}}_{\hat{b}} \in \mathbb{R}^{d + 1} $,
        $ \nabla \Vert\mathbf{y} - \mathbf{X}_\mathbf{1}
        \hat{\mathbf{w}}_{\hat{b}}\Vert_2^2 = \mathbf{0} $. See that
        \begin{equation*}
            \nabla\Vert\mathbf{y} -
            \mathbf{X}_\mathbf{1}\hat{\mathbf{w}}_{\hat{b}}\Vert_2^2 =
            -2\mathbf{X}_\mathbf{1}^\top(\mathbf{y} -
            \mathbf{X}_\mathbf{1}\hat{\mathbf{w}}_{\hat{b}}) = \mathbf{0}
            \Leftrightarrow \mathbf{X}_\mathbf{1}^\top\mathbf{y} =
            \mathbf{X}_\mathbf{1}^\top\mathbf{X}_\mathbf{1}
            \hat{\mathbf{w}}_{\hat{b}}
        \end{equation*}

        \item
        $ \mathbf{X}_\mathbf{1}^\top\mathbf{X}_\mathbf{1} \in
        \mathbb{R}^{(d + 1) \times (d + 1)} $ is full rank\footnote{
            See first answer of this \ahref{%
                https://math.stackexchange.com/questions/349738/prove-%
                operatornamerankata-operatornameranka-for-any-a-in-m-m-times-n%
            }{math.stackexchange} question for proof.
        }, i.e. invertible\footnote{
            In practice this inverse is rarely computed. For example,
            \ahref{%
                https://docs.scipy.org/doc/scipy/reference/generated/%
                scipy.linalg.lstsq.html%
            }{\texttt{scipy.linalg.lstsq}} uses LAPACK's \ahref{%
                http://www.netlib.org/lapack/explore-html/db/d6a/%
                dgelsd_8f.html%
            }{\texttt{gelsd}} by default, which uses the singular value
            decomposition.
        },
        and therefore
        \begin{equation} \label{lr_ols_aug}
            \hat{\mathbf{w}}_{\hat{b}} = \big(
                \mathbf{X}_\mathbf{1}^\top\mathbf{X}_\mathbf{1}
            \big)^{-1}\mathbf{X}_\mathbf{1}^\top\mathbf{y}
        \end{equation}
        Or, written in terms of the original $ \mathbf{X}, \hat{\mathbf{w}},
        \hat{b} $,
        \begin{equation} \label{lr_ols_block}
            \begin{bmatrix}
                \ \hat{\mathbf{w}} \ \\ \ \hat{b} \
            \end{bmatrix} =
            \left(
                \begin{bmatrix}
                    \ \mathbf{X}^\top \ \\ \ \mathbf{1}^\top \
                \end{bmatrix}
                \begin{bmatrix}
                    \ \mathbf{X} & \mathbf{1} \
                \end{bmatrix}
            \right)^{-1}
            \begin{bmatrix}
                \ \mathbf{X} & \mathbf{1} \
            \end{bmatrix}^\top\mathbf{y}
        \end{equation}
    \end{itemize}
\end{frame}

\subsection{On centered data}

%\begin{frame}{On centered data}
%    \begin{itemize}
%        \item
%        Consider: can we estimate $ \mathbf{w}, b $ separately?
%
%        \item
%        \textit{Proposition.} Let $ X : \Omega \rightarrow \mathbb{R}^d $,
%        $ Y : \Omega \rightarrow \mathbb{R} $ be random variables on
%        $ (\Omega, \mathcal{F}, \mathbb{P}) $, and let $ \varepsilon \sim
%        \mathcal{N}(0, \sigma^2) $ be independent of $ X $, $ \sigma \in
%        (0, \infty) $. For $ \mathbf{w} \in \mathbb{R}^d $,
%        $ b \in \mathbb{R} $,
%        $ Y = \mathbf{w}^\top X + b + \varepsilon \Rightarrow b = \mu_Y -
%        \mathbf{w}^\top\mu_X $.
%
%        \item
%        \textit{Proof.} Taking expectations, $ \mu_Y = \mathbf{w}^\top\mu_X +
%        b \Leftrightarrow b = \mu_Y - \mathbf{w}^\top\mu_X $.
%
%        \item
%        \textit{Remark.} For any model $ Y = f(X) + b + \varepsilon $,
%        $ b = \mu_Y - \mu_{f(X)} $.
%
%        \item
%        \textit{Corollary.} $ \mu_X = \mathbf{0} \Rightarrow b = \mu_Y $.
%
%        \item
%        This implies that $ \mathbf{w} $
%        
%    \end{itemize}        
%\end{frame}

\begin{frame}{On centered data}
    \begin{itemize}
        \item
        Consider: can we estimate $ \mathbf{w}, b $ separately\footnote{
            For $ h : X(\Omega) \rightarrow \mathbb{R} $, any model
            $ Y = h(X) + b + \varepsilon \Rightarrow b = \mu_Y - \mu_{h(X)} $.
        }?

        \item
        Let $ \mathbf{X} \in \mathbb{R}^{N \times d} $ be the input matrix,
        $ \mathbf{y} \in \mathbb{R}^N $ be the response vector. Let
        $ \bar{\mathbf{x}} \triangleq \frac{1}{N}\mathbf{X}^\top\mathbf{1} $,
        $ \bar{y} \triangleq \frac{1}{N}\mathbf{y}^\top\mathbf{1} $,
        $ \tilde{\mathbf{X}} \triangleq \mathbf{X} -
        \mathbf{1}\bar{\mathbf{x}}^\top $, $ \tilde{\mathbf{y}} \triangleq
        \mathbf{y} - \bar{y}\mathbf{1} $. Consider
        \begin{equation*}
            \begin{array}{lll}
                \displaystyle\min_{\mathbf{w}, b} &
                \Vert\mathbf{y} - \mathbf{Xw} - b\mathbf{1}\Vert_2^2 =
                \displaystyle\min_{\mathbf{w}, b} & g(\mathbf{w}, b)
            \end{array}
        \end{equation*}
        This is (\aref{lr_obj_aug}), but using $ \mathbf{X} $ and writing
        $ \mathbf{w}, b $ separately. $ g $ is convex in $ \mathbf{w}, b
        \Rightarrow \nabla_\mathbf{w}g(\hat{\mathbf{w}}, \hat{b}) =
        \mathbf{0}, \frac{\partial g}{\partial b}(\hat{\mathbf{w}}, \hat{b})
        = 0 $\footnote{
            Since $ f $ convex, these are necessary and sufficient conditions
            for optimality \cite{bv_convex_opt}.
        }. See that
        \begin{equation*}
            \nabla_\mathbf{w} g(\hat{\mathbf{w}}, \hat{b}) =
            -2\mathbf{X}^\top(\mathbf{y} - \mathbf{X}\hat{\mathbf{w}} -
            \hat{b}\mathbf{1}) = \mathbf{0} \Leftrightarrow
            \mathbf{X}^\top\mathbf{X}\hat{\mathbf{w}} =
            \mathbf{X}^\top\mathbf{y} - \hat{b}\mathbf{X}^\top\mathbf{1}            
        \end{equation*}
        \begin{equation*}
            \begin{split}
                \frac{\partial g}{\partial b}(\hat{\mathbf{w}}, \hat{b}) =
	            -2\mathbf{1}^\top(\mathbf{y} - \mathbf{X}\hat{\mathbf{w}} -
	            \hat{b}\mathbf{1}) = 0 & \Leftrightarrow
	            \hat{b}\mathbf{1}^\top\mathbf{1} = \mathbf{y}^\top\mathbf{1} -
	            \hat{\mathbf{w}}^\top\mathbf{X}^\top\mathbf{1} \\
	            & \Rightarrow b = \bar{y} - \hat{\mathbf{w}}^\top\bar{\mathbf{x}}
            \end{split}
        \end{equation*}        
    \end{itemize}

    % spacing for footnote
    \medskip
\end{frame}

\begin{frame}{On centered data}
    \begin{itemize}
        \item
        Note that for $ \mathbf{1} \in \mathbb{R}^N,
        \mathbf{1}^\top\mathbf{1} = N $. Substituting,
        \begin{equation*}
            \begin{split}
	            & \mathbf{X}^\top\mathbf{X}\hat{\mathbf{w}} =
	            \mathbf{X}^\top\mathbf{y} - \bar{y}\mathbf{X}^\top\mathbf{1} +
	            \mathbf{X}^\top\mathbf{1}\bar{\mathbf{x}}^\top\hat{\mathbf{w}} \\
	            \Rightarrow \ & \mathbf{X}^\top
	            (\mathbf{X} - \mathbf{1}\bar{\mathbf{x}}^\top)\hat{\mathbf{w}} =
	            \mathbf{X}^\top(\mathbf{y} - \bar{y}\mathbf{1}) \\
	            \Rightarrow \ & \mathbf{X}^\top\left(
	                \mathbf{X} - \frac{1}{N}\mathbf{11}^\top\mathbf{X}
	            \right)\hat{\mathbf{w}} = \mathbf{X}^\top\left(
	                \mathbf{y} - \frac{1}{N}\mathbf{11}^\top\mathbf{y}
	            \right) \\
	            \Rightarrow \ & \mathbf{X}^\top\left(
	                \mathbf{I} - \frac{1}{N}\mathbf{11}^\top
	            \right)\mathbf{X}\hat{\mathbf{w}} = \mathbf{X}^\top\left(
	                \mathbf{I} - \frac{1}{N}\mathbf{11}^\top
	            \right)\mathbf{y}
            \end{split}
        \end{equation*}

        \item
        Recall that $ \mathbf{I} - \frac{1}{N}\mathbf{11}^\top \triangleq
        \tilde{\mathbf{C}} $ is the \textit{centering matrix}. Note
        $ \tilde{\mathbf{C}} = \tilde{\mathbf{C}}^\top $ (symmetric),
        $ \tilde{\mathbf{C}}\tilde{\mathbf{C}} = \tilde{\mathbf{C}} $
        (idempotent), so the last equation becomes
        \begin{equation*}
            \big(\tilde{\mathbf{C}}\mathbf{X}\big)^\top\big(
                \tilde{\mathbf{C}}\mathbf{X}
            \big)\hat{\mathbf{w}} =
            \big(\tilde{\mathbf{C}}\mathbf{X}\big)^\top\mathbf{y}
        \end{equation*}
    \end{itemize}
\end{frame}

\begin{frame}{On centered data}
    \begin{itemize}
        \item
        Note that the centering matrix $ \tilde{\mathbf{X}} \triangleq
        \mathbf{X} - \mathbf{1}\bar{\mathbf{x}}^\top =
        \tilde{\mathbf{C}}\mathbf{X} $. Therefore,
        \begin{equation} \label{lr_ols_cent}
            \begin{split}
                \hat{\mathbf{w}} & = \big(
	                \tilde{\mathbf{X}}^\top\tilde{\mathbf{X}}
	            \big)^{-1}\tilde{\mathbf{X}}^\top\mathbf{y} \\
	            \hat{b} & = \bar{y} - \hat{\mathbf{w}}^\top\bar{\mathbf{x}}
            \end{split}
        \end{equation}
        Note that $ \tilde{\mathbf{X}}^\top\tilde{\mathbf{X}} $ invertible
        $ \Leftrightarrow \operatorname{rk}\big(\tilde{\mathbf{X}}\big) = d $.

        \item
        (\aref{lr_ols_cent}) says that by using $ \tilde{\mathbf{X}} $, we can
        compute $ \hat{\mathbf{w}} $ independently of $ \hat{b} $.

        \item
        \textit{Remark.} If $ \mathbf{X} = \tilde{\mathbf{X}} $, i.e. the
        inputs are already centered, then $ \hat{b} = \bar{y} $. If
        $ \mathbf{X} = \tilde{\mathbf{X}}, \mathbf{y} = \tilde{\mathbf{y}} $,
        i.e. both inputs and outputs are centered, then $ \hat{b} = 0 $.

        \item
        \alert{No intercept is needed if both inputs and outputs are centered.}
    \end{itemize}
\end{frame}

\section{Generalizations}

\subsection{Weighted least squares}

\begin{frame}{Weighted least squares}
    \begin{itemize}
        \item
        Let $ \mathbf{\Gamma} \triangleq \operatorname{diag}(\gamma_1, \ldots
        \gamma_N)$ , $ \gamma_1, \ldots \gamma_N > 0 $. Consider
        \begin{equation} \label{wlr_obj}
            \begin{array}{ll}
                \displaystyle\min_{\mathbf{w}, b} &
                \left\Vert\mathbf{\Gamma}^{1 / 2}(\mathbf{y} - \mathbf{Xw} - 
                b\mathbf{1})\right\Vert_2^2
            \end{array}
        \end{equation}
        $ \mathbf{\Gamma} $ weights the components of the squared residuals.
        The solution $ \hat{\mathbf{w}}_\mathbf{\Gamma},
        \hat{b}_\mathbf{\Gamma} $ of (\aref{wlr_obj}), is given by the
        equations\footnote{
            We derive $ \hat{\mathbf{w}}_\mathbf{\Gamma},
            \hat{b}_\mathbf{\Gamma} $ using similar steps as those used to
            derive $ \hat{\mathbf{w}}, \hat{b} $ in (\aref{lr_ols_cent}), but
            $ \tilde{\mathbf{C}} $ is replaced with
            $ \tilde{\mathbf{C}}_\mathbf{\Gamma} $ and $ \bar{\mathbf{x}},
            \bar{y} $ are replaced with
            $ \bar{\mathbf{x}}_\mathbf{\Gamma}, \bar{y}_\mathbf{\Gamma} $.
        }
        \begin{equation} \label{wlr_sol}
            \begin{split}
                \hat{\mathbf{w}}_\mathbf{\Gamma} & = \big(
                    \tilde{\mathbf{X}}_\mathbf{\Gamma}^\top
                    \mathbf{\Gamma}\tilde{\mathbf{X}}_\mathbf{\Gamma}
                \big)^{-1}\tilde{\mathbf{X}}_\mathbf{\Gamma}^\top
                \mathbf{\Gamma}\mathbf{y}_\mathbf{\Gamma} \\
                \hat{b}_\mathbf{\Gamma} & = \bar{y}_\mathbf{\Gamma} -
                \hat{\mathbf{w}}_\mathbf{\Gamma}^\top
                \bar{\mathbf{x}}_\mathbf{\Gamma}
            \end{split}
        \end{equation}

        \item
        $ \bar{\mathbf{x}}_\mathbf{\Gamma} \triangleq
        \frac{1}{\mathbf{1}^\top\mathbf{\Gamma 1}}\mathbf{X}^\top
        \mathbf{\Gamma 1}, \bar{y}_\mathbf{\Gamma} \triangleq
        \frac{1}{\mathbf{1}^\top\mathbf{\Gamma 1}}\mathbf{y}^\top
        \mathbf{\Gamma 1} $ are \textit{weighted averages}.
        $ \tilde{\mathbf{X}}_\mathbf{\Gamma} \triangleq
        \tilde{\mathbf{C}}_\mathbf{\Gamma}\mathbf{X} $, where
        $ \tilde{\mathbf{C}}_\mathbf{\Gamma} \triangleq \mathbf{I} -
        \frac{1}{\mathbf{1}^\top\mathbf{\Gamma 1}}
        \mathbf{11}^\top\mathbf{\Gamma} $ is the
        \textit{weighted centering matrix}.
    \end{itemize}
\end{frame}

\subsection{Ridge regression}

\begin{frame}{Ridge regression}
    \begin{itemize}
        \item
        At its simplest, ridge regression serves as a way to deflate, or
        \textit{regularize}, the coefficients of $ \hat{\mathbf{w}} $. For
        $ \lambda \in [0, \infty) $, consider
        \begin{equation} \label{ridge_obj}
            \begin{array}{ll}
                \displaystyle\min_{\mathbf{w}, b} &
                \Vert\mathbf{y} - \mathbf{Xw} - b\mathbf{1}\Vert_2^2 +
                \lambda\Vert\mathbf{w}\Vert_2^2
            \end{array}
        \end{equation}
        $ \lambda\Vert\mathbf{w}\Vert_2^2 $ penalizes a candidate point
        $ \mathbf{w} $ based on its $ \ell^2 $-norm. $ \lambda $ is a free
        parameter that controls the penalization tradeoff.

        \item
        The solution $ \hat{\mathbf{w}}_\lambda, \hat{b}_\lambda $ of
        (\aref{ridge_obj}) is given by\footnote{
            We derive $ \hat{\mathbf{w}}_\lambda, \hat{b}_\lambda $ using
            almost the exact same steps used to derive $ \hat{\mathbf{w}},
            \hat{b} $.
        }
        \begin{equation} \label{ridge_sol}
            \begin{split}
                \hat{\mathbf{w}}_\lambda & = \big(
	                \tilde{\mathbf{X}}^\top\tilde{\mathbf{X}} +
	                \lambda\mathbf{I}
	            \big)^{-1}\tilde{\mathbf{X}}^\top\mathbf{y} \\
	            \hat{b}_\lambda & = \bar{y} - \hat{\mathbf{w}}_\lambda^\top
	            \bar{\mathbf{x}}
            \end{split}
        \end{equation}
    \end{itemize}
\end{frame}

\begin{frame}{Ridge regression}
    \begin{itemize}
        \item
        Why use $ \lambda\Vert\mathbf{w}\Vert_2^2 $ as the
        regularizer\footnote{
            The list is not exhaustive. For example,
            $ \hat{\mathbf{w}}_\lambda $ has lower variance than
            $ \hat{\mathbf{w}} $, and
            $ \tilde{\mathbf{X}}^\top\tilde{\mathbf{X}} +
            \lambda\mathbf{I} $ is well-conditioned even if
            $ \tilde{\mathbf{X}}^\top\tilde{\mathbf{X}} $ itself is poorly
            conditioned.
        }?
        \begin{enumerate}
            \item
            Analytically tractable, i.e. closed form solution available.

            \item
            Interpretable as a Bayesian prior over the coefficients.
        \end{enumerate}

        \item
        \alert{
            $ \hat{\mathbf{w}}_\lambda $ is affected by the
            \textit{scaling} of the input matrix $ \mathbf{X} $. Why?
        }

        \item
        Suppose the sample covariance matrix of $ \mathbf{X} $,
        $ \mathbf{C}_\mathbf{X} $, is such that
        $ \mathbf{C}_\mathbf{X} = \mathbf{I} $, i.e. uncorrelated, unit
        variance features. $ \tilde{\mathbf{X}} $ has orthogonal columns.

        \item
        Since $ \mathbf{C}_\mathbf{X} \triangleq
        \frac{1}{N}\tilde{\mathbf{X}}^\top\tilde{\mathbf{X}} = \mathbf{I} $,
        then by (\aref{lr_ols_cent}), the OLS solution is
        \begin{equation} \label{lr_ols_ortho}
            \begin{split}
	            \hat{\mathbf{w}} & = (N\mathbf{I})^{-1}\tilde{\mathbf{X}}^\top
	            \mathbf{y} = \frac{1}{N}\tilde{\mathbf{X}}^\top\mathbf{y} \\
	            \hat{b} & = \bar{y} - \hat{\mathbf{w}}^\top\bar{\mathbf{x}}
	        \end{split}
        \end{equation}
    \end{itemize}
\end{frame}

\begin{frame}{Ridge regression}
    \begin{itemize}
        \item
        By (\aref{ridge_sol}), (\aref{lr_ols_ortho}), the ridge solution
        $ \hat{\mathbf{w}}_\lambda $ is
        \begin{equation} \label{ridge_sol_ortho}
            \begin{split}
	            \hat{\mathbf{w}}_\lambda & =
	            (N\mathbf{I} + \lambda\mathbf{I})^{-1}
	            \tilde{\mathbf{X}}^\top\mathbf{y} =
	            \frac{1}{N + \lambda}\tilde{\mathbf{X}}^\top\mathbf{y} =
	            \boxed{\frac{N}{N + \lambda}\hat{\mathbf{w}}} \\
	            \hat{b}_\lambda & = \bar{y} - \hat{\mathbf{w}}_\lambda^\top
	            \bar{\mathbf{x}} = \bar{y} - \frac{N}{N + \lambda}
	            (\bar{y} - \hat{b}) =
	            \boxed{
	                \frac{N}{N + \lambda}\hat{b} +
	                \frac{\lambda}{N + \lambda}\bar{y}
	            }
            \end{split}
        \end{equation}

        \item
        (\aref{ridge_sol_ortho}) tells us that in the case of uncorrelated,
        unit variance input features, $ \hat{\mathbf{w}}_\lambda $
        deflates $ \hat{\mathbf{w}} $ by $ N / (N + \lambda) $, and
        $ \hat{b}_\lambda $ is a convex combination of $ \hat{b} $ and
        $ \bar{y} $. Clearly $ \lim_{\lambda \rightarrow \infty}
        \hat{\mathbf{w}}_\lambda = \mathbf{0} $,
        $ \lim_{\lambda \rightarrow \infty}\hat{b}_\lambda = \bar{y} $.

        \item
        Now suppose that $ \mathbf{C}_\mathbf{X} \triangleq
        \operatorname{diag}(\nu_1, \ldots \nu_d) $,
        $ \nu_1, \ldots \nu_d > 0 $. By (\aref{lr_ols_cent}),
        \begin{equation} \label{lr_ols_diag}
            \begin{split}
                \hat{\mathbf{w}}_\nu & = (N\mathbf{C}_\mathbf{X})^{-1}
	            \tilde{\mathbf{X}}^\top\mathbf{y} =
	            \frac{1}{N}\operatorname{diag}\left(
	                \nu_1^{-1}, \ldots \nu_d^{-1}
	            \right)
	            \tilde{\mathbf{X}}^\top\mathbf{y} \\
	            \hat{b}_\nu & = \bar{y} - \hat{\mathbf{w}}_\nu^\top
	            \bar{\mathbf{x}}
            \end{split}
        \end{equation}
    \end{itemize}
\end{frame}

\begin{frame}{Ridge regression}
    \begin{itemize}
        \item
        Similarly, by (\aref{ridge_sol}), the new ridge solution
        $ \hat{\mathbf{w}}_{\lambda, \nu}, \hat{b}_\nu $ is
        \begin{equation} \label{ridge_sol_diag}
            \begin{split}
                \hat{\mathbf{w}}_{\lambda, \nu} & =
                (N\mathbf{C}_\mathbf{X} + \lambda\mathbf{I})^{-1}
                \tilde{\mathbf{X}}^\top\mathbf{y} \\
                & = \operatorname{diag}\left(
                    \frac{1}{N\nu_1 + \lambda}, \ldots
                    \frac{1}{N\nu_d + \lambda}
                \right)\tilde{\mathbf{X}}^\top\mathbf{y} \\
                & = \boxed{
                    \operatorname{diag}\left(
	                    \frac{N\nu_1}{N\nu_1 + \lambda}, \ldots
	                    \frac{N\nu_d}{N\nu_d + \lambda}
	                \right)\hat{\mathbf{w}}_\nu
	            } \\
	            \hat{b}_{\lambda, \nu} & = \bar{y} -
                \hat{\mathbf{w}}_{\lambda, \nu}^\top\bar{x} \\
                & = \boxed{
	                \bar{y} - \hat{\mathbf{w}}_\nu^\top
	                \operatorname{diag}\left(
		                \frac{N\nu_1}{N\nu_1 + \lambda}, \ldots
	                    \frac{N\nu_d}{N\nu_d + \lambda}
	                \right)\bar{\mathbf{x}}
	            }
            \end{split}
        \end{equation}
    \end{itemize}
\end{frame}

\begin{frame}{Ridge regression}
    \begin{itemize}
        \item 
        (\aref{ridge_sol_diag}) shows that the shrinkage applied by
        $ \hat{\mathbf{w}}_{\lambda, \nu} $ to the new OLS coefficients
        $ \hat{\mathbf{w}}_\nu $ now depends on the feature variances
        $ \nu_1, \ldots \nu_d $.

        \item
        $ i $th element $ \hat{w}_{\nu_i} $ of $ \hat{\mathbf{w}}_\nu $ is
        deflated by $ N\nu_i / (N\nu_i + \lambda) $. For fixed $ \lambda $,
        \begin{equation*}
            \begin{array}{cc}
                \displaystyle\lim_{\nu_i \rightarrow 0^+}
                \frac{N\nu_i}{N\nu_i + \lambda} = 0 \quad & \quad
                \displaystyle\lim_{\nu_i \rightarrow \infty}
                \frac{N\nu_i}{N\nu_i + \lambda} = 1
            \end{array}
        \end{equation*}

        \item
        In other words, given a fixed choice of $ \lambda $, if
        $ \nu_i\rightarrow 0 $, then $ \hat{w}_{\nu_i} $ is deflated
        \textbf{more}, while if $ \nu_i \rightarrow \infty $, then
        $ \hat{w}_{\nu_i} $ is deflated \textbf{less}.

        \item
        \alert{
            Ridge regression implicitly assumes that higher-variance features
            are \textbf{more predictive} and therefore shrinks the relevant
            coefficients \textbf{less}
        }.

        \item
        The assumption may be helpful or harmful, depending on context.
    \end{itemize}
\end{frame}

% BibTeX slide for references. should use either acm or ieeetr style
\begin{frame}{References}
    \bibliographystyle{acm}
    % relative path may need to be updated depending on .tex file location
    \bibliography{../master_bib}
\end{frame}

\end{document}